\documentclass[11pt]{article}

\usepackage[margin=1in]{geometry}
\usepackage{amsmath,amssymb,amsthm,mathtools}
\usepackage[colorlinks=true,linkcolor=black,urlcolor=blue,citecolor=blue]{hyperref}

\newcommand{\studentname}{YOUR NAME}
\newcommand{\workingtitle}{A precise working title}

\title{\workingtitle\\[0.4em]
\large Research-question note for CMPUT 654}
\author{\studentname}
\date{Fall 2026}

\begin{document}
\maketitle

\section{The question}
State the question as precisely as possible. Define the mathematical objects,
learner, information protocol, resource constraint, or evaluation criterion
needed to make it meaningful.

\section{Why this question matters}
Explain what uncertainty about modern AI the answer would clarify. Distinguish
the mathematical question from broader empirical or philosophical claims.

\section{What is known}
Summarize the most relevant positive results, negative results, constructions,
or empirical evidence. State their assumptions and cite the sources.

\section{What remains unknown}
Identify the gap. Explain which assumptions may be essential and what kind of
result, counterexample, or experiment would count as progress.

\section{Revision note}
For the final submission, briefly explain how feedback changed or narrowed the
question.

\begin{thebibliography}{9}
\bibitem{placeholder}
Author names.
\newblock \emph{Paper title}.
\newblock Venue or arXiv identifier, year.
\end{thebibliography}

\end{document}
